% W11 Dynamics of wheeled robots -- booklet v3 (condensed). Numbers: vcheck/v2_w11.py.
% Slide notation: {}^C F_G, {}^C N_G, {}^G I, {}^C P_G, u, v, thetadot.  Deck = Newton--Euler of a planar chassis (no control content).
\newcommand{\wEz}{\boldsymbol\zeta}
\newcommand{\wEx}{\boldsymbol\xi}
\newcommand{\wEQC}{{}^{C}\mathbf Q}
\newcommand{\wEQO}{{}^{O}\mathbf Q}
\newcommand{\wEPG}{{}^{C}\mathbf P_G}
\newcommand{\wEth}{\dot\theta}
\newcommand{\wEthh}{\ddot\theta}
\newcommand{\wEI}{{}^{G}I}

\week{W11}{Dynamics of wheeled robots}{11}
\begin{multicols}{3}

%==================================================================
\hd{Frames, unknowns, the two problems \pg{9--12}}\nopagebreak[4]
\slidefig[0.72\linewidth]{w11-lec}{12}{29 119 586 194}{$\{O\}$ inertial, $\{C\}$ on chassis, $\{G\}$ at COM, same orientation as $\{C\}$}
\begin{fbx}[Setup]
$\wEx=(x,y,\theta)$ pose of $C$ in $\{O\}$; $\wEz=(u,v,\wEth)$ chassis velocity in $\{C\}$; $\dot{\wEx}=\mathbf J\wEz$, $\mathbf J={}^O\mathbf R_C=\begin{bmatrix}c\theta&-s\theta&0\\s\theta&c\theta&0\\0&0&1\end{bmatrix}$.
$\mathbf Q=(f_x,f_y,\tau_z)$ net load applied at $C$ (wheels+ground), $\wEQC$ or $\wEQO$: always name the frame.
$m$; $\wEI={}^GI_{zz}$ about $G$; $\wEPG=({}^Cx_G,{}^Cy_G,0)$ constant (both frames turn with the robot); $C=G$ simplifies everything.
\end{fbx}
\premise[forward \pg{10}]{$\mathbf Q$, $m$, $\wEI$, geometry}{$(\wEx,\dot{\wEx},\ddot{\wEx})$: simulate. \textbf{Inverse \pg{11}:} given the motion, find $\mathbf Q$ (control)}
No control in W11 (it is W8--W9). Do not mix $\wEx$ (pose, $\{O\}$) with $\wEz$ (velocity, $\{C\}$).

%==================================================================
\hd{Velocity and acceleration of $G$ \pg{18--24}}\nopagebreak[4]
\begin{deriv}[velocity \pg{19--20}]
${}^O\mathbf P_G={}^O\mathbf P_C+{}^O\mathbf R_C\,\wEPG\Rightarrow{}^O\dot{\mathbf P}_G={}^O\dot{\mathbf P}_C+{}^O\boldsymbol\omega\times{}^O\mathbf R_C\wEPG$ (${}^C\dot{\mathbf P}_G=\mathbf 0$, $G$ fixed in chassis; $\dot{\mathbf R}\mathbf p=\boldsymbol\omega\times\mathbf R\mathbf p$). Rotations keep cross products, so factor ${}^O\mathbf R_C$:
\end{deriv}
\begin{fbx}[Velocity of $G$ \pg{19--20}]
\[{}^O\dot{\mathbf P}_G={}^O\mathbf R_C\Big[{}^C\dot{\mathbf P}_C+{}^C\boldsymbol\omega_C\times\wEPG\Big],\ \ {}^C\dot{\mathbf P}_C=\begin{bmatrix}u\\v\\0\end{bmatrix},\ {}^C\boldsymbol\omega_C=\begin{bmatrix}0\\0\\\wEth\end{bmatrix}\]
\[\Rightarrow\ {}^C\dot{\mathbf P}_G=\begin{bmatrix}u\\v\\0\end{bmatrix}+\begin{bmatrix}-\wEth\,{}^Cy_G\\\wEth\,{}^Cx_G\\0\end{bmatrix}=\begin{bmatrix}u-\wEth\,{}^Cy_G\\v+\wEth\,{}^Cx_G\\0\end{bmatrix}\]
$(0,0,k)\times(a,b,0)=(-kb,\,ka,\,0)$: used on every line of this week.
\end{fbx}
\begin{worked}[velocity, ours]
$G=(0.3,0.2)$, $u=1.0$, $v=0.5$, $\wEth=2$: $\{C\}$: $(1.0-0.4,\ 0.5+0.6)=(\mathbf{0.6},\mathbf{1.1})$ m/s. At $\theta=30^\circ$: $\mathbf J(0.6,1.1)=(\mathbf{-0.0304},\mathbf{1.2526})$; speed $1.253$ in both. \chk
\end{worked}
\begin{deriv}[acceleration \pg{22}]
Differentiate again (product rule); the term $\boldsymbol\omega\times{}^C\dot{\mathbf P}_G=0$. With $\dot{\mathbf R}\mathbf p=\mathbf R(\boldsymbol\omega\times\mathbf p)$ factor ${}^O\mathbf R_C$ and premultiply by its inverse:
\end{deriv}
\begin{fbx}[Acceleration of $G$ in $\{C\}$: four terms \pg{22--24}]
\[\boxed{{}^C\ddot{\mathbf P}_G={}^C\ddot{\mathbf P}_C+{}^C\dot{\boldsymbol\omega}_C\times\wEPG+{}^C\boldsymbol\omega_C\times{}^C\dot{\mathbf P}_C+{}^C\boldsymbol\omega_C\times{}^C\boldsymbol\omega_C\times\wEPG}\]
${}^C\ddot{\mathbf P}_C=(\dot u,\dot v,0)$, ${}^C\dot{\boldsymbol\omega}_C=(0,0,\wEthh)$. Using $[\boldsymbol\omega\times]=\begin{bmatrix}0&-\wEth&0\\\wEth&0&0\\0&0&0\end{bmatrix}$ ($\wEthh$ for $\dot{\boldsymbol\omega}$):
\[\underbrace{\begin{bmatrix}-y_G\wEthh\\x_G\wEthh\\0\end{bmatrix}}_{\text{tangential}}+\underbrace{\begin{bmatrix}-v\wEth\\u\wEth\\0\end{bmatrix}}_{\text{Coriolis-like}}+\underbrace{\begin{bmatrix}-x_G\wEth^2\\-y_G\wEth^2\\0\end{bmatrix}}_{\text{centripetal}}\!,\qquad
\begin{aligned}{}^C\ddot x_G&=\dot u-y_G\wEthh-v\wEth-x_G\wEth^2\\{}^C\ddot y_G&=\dot v+x_G\wEthh+\boxed{u}\wEth-y_G\wEth^2\end{aligned}\]
\end{fbx}
\begin{slideerr}[pp.\,24--26: $v\dot\theta$ should be $u\dot\theta$ in $y$]
Slides print $+v\dot\theta$ in ${}^C\ddot y_G$, ${}^CF_y$, ${}^Cf_y$. Correct: $+u\dot\theta$ (the slide's own cross product gives it; p.\,26 matrix uses $u$; constant $u$ on a circle needs $a_y=u\dot\theta$). The $x$ row is right. \chk
\end{slideerr}
\begin{worked}[all four terms, ours]
$u{=}1.5$, $v{=}0.2$, $\wEth{=}0.8$, $\dot u{=}0.3$, $\dot v{=}{-}0.1$, $\wEthh{=}0.5$, $G{=}(0.2,0.1)$: $\ddot{\mathbf P}_C(0.3,-0.1)$ + tangential $(-0.05,0.10)$ + Coriolis-like $(-0.16,1.20)$ + centripetal $(-0.128,-0.064)$ $=\mathbf a_G=(\mathbf{-0.038},\mathbf{1.136})$ m/s$^2$ (printed $+v\dot\theta$: $a_y{=}0.096$, wrong). \chk
\end{worked}

%==================================================================
\hd{Newton--Euler in $\{C\}$: the equation of motion \pg{14--26}}\nopagebreak[4]
\begin{fbx}[Balance \pg{15--17}]
Inertial force/torque at $G$: ${}^C\mathbf F_G=m\,{}^C\ddot{\mathbf P}_G$, ${}^CN_z=\wEI\wEthh$ ($\boldsymbol\omega\times\wEI\boldsymbol\omega=\mathbf 0$ in the plane). Applied $({}^Cf_x,{}^Cf_y)$, ${}^Cn_z$ act at $C$; lever arm $G\to C$ is $-\wEPG$:
\[{}^CF_x={}^Cf_x,\ {}^CF_y={}^Cf_y,\quad{}^CN_z={}^Cn_z+\underbrace{(-\wEPG)\times{}^C\mathbf f}_{=\,y_Gf_x-x_Gf_y}\]
\end{fbx}
\begin{fbx}[Newton--Euler equations \pg{25}]
\[{}^Cf_x=m(\dot u-y_G\wEthh-v\wEth-x_G\wEth^2),\quad{}^Cf_y=m(\dot v+x_G\wEthh+u\wEth-y_G\wEth^2)\]
\[{}^Cn_z=\wEI\wEthh-y_G\,{}^Cf_x+x_G\,{}^Cf_y\]
\end{fbx}
\begin{deriv}[to matrix form \added]
Insert $f_x,f_y$ in $n_z$; the $x_Gy_G\wEth^2$ terms cancel: ${}^Cn_z=-my_G\dot u+mx_G\dot v+\big(\wEI+m(x_G^2+y_G^2)\big)\wEthh+m\wEth(x_Gu+y_Gv)$.
\end{deriv}
\begin{fbx}[Matrix form $\mathbf M\dot{\wEz}+c(\wEz)=\wEQC$ \pg{26}]
\[\begin{bmatrix}m&0&-my_G\\0&m&mx_G\\-my_G&mx_G&\wEI{+}m(x_G^2{+}y_G^2)\end{bmatrix}\!\begin{bmatrix}\dot u\\\dot v\\\wEthh\end{bmatrix}+\begin{bmatrix}-m\wEth(v+x_G\wEth)\\m\wEth(u-y_G\wEth)\\m\wEth(x_Gu+y_Gv)\end{bmatrix}=\begin{bmatrix}{}^Cf_x\\{}^Cf_y\\{}^Cn_z\end{bmatrix}\]
$M_{33}$ = inertia about $C$ (parallel axis). Simulate: $\dot{\wEz}=\mathbf M^{-1}(\wEQC-c)$.
\textbf{$C=G$:} $\mathbf M=\operatorname{diag}(m,m,\wEI)$, $c=(-mv\wEth,\ mu\wEth,\ 0)$: $m(\dot u-v\wEth)=f_x$, $m(\dot v+u\wEth)=f_y$, $\wEI\wEthh=n_z$.
\end{fbx}
\begin{recipe}[Derive the EoM (W13 Q4 style)]
\begin{enumerate}
\item Draw $\{O\},\{C\},\{G\}$, loads at $C$, inertial loads at $G$, $(x_G,y_G)$.
\item ${}^C\dot{\mathbf P}_G=(u-\wEth y_G,\ v+\wEth x_G)$.
\item $\mathbf a_G$ from the four terms (expand each cross product).
\item $f_x=ma_x$, $f_y=ma_y$, $n_z=\wEI\wEthh-(y_Gf_x-x_Gf_y)$.
\item Collect $\to\mathbf M\dot{\wEz}+c=\wEQC$.
\item \textbf{Check:} $\mathbf M$ symmetric; $c=0$ at rest; $C=G$ gives diag; units N, N\,m.
\end{enumerate}
\end{recipe}
\textbf{Lagrange alternative \pg{27}, \added:} $T=\tfrac12m|\dot{\mathbf P}_G|^2+\tfrac12\wEI\wEth^2=\tfrac12\wEz^\top\mathbf M\wEz$, $\tfrac{d}{dt}\tfrac{\partial T}{\partial\dot{\wEx}}-\tfrac{\partial T}{\partial\wEx}=\wEQO$: same motion, no cross products \chk
\begin{worked}[offset COM, at rest]
$m{=}8$, $\wEI{=}0.3$, $G{=}(0.05,0.02)$, ${}^C\mathbf f{=}(6,0)$, ${}^Cn_z{=}0.5$, $c{=}\mathbf 0$; $M_{33}=0.3+8(0.0025+0.0004)=0.3232$:
\[\begin{bmatrix}8&0&-0.16\\0&8&0.4\\-0.16&0.4&0.3232\end{bmatrix}\!\begin{bmatrix}\dot u\\\dot v\\\wEthh\end{bmatrix}=\begin{bmatrix}6\\0\\0.5\end{bmatrix}\Rightarrow\wEthh=\mathbf{2.067},\ \dot v=\mathbf{-0.103},\ \dot u=\mathbf{0.791}\]
($M_{33}=\wEI$ alone gives $2.24$: wrong.) \chk
\end{worked}

%==================================================================
\hd{From $\{C\}$ to $\{O\}$ \pg{28--34}}\nopagebreak[4]
\begin{fbx}[Kinematic link \pg{30}]
\[\dot{\wEx}=\mathbf J\wEz,\qquad \dot{\wEz}=\mathbf J^{-1}\big[\ddot{\wEx}-\dot{\mathbf J}\mathbf J^{-1}\dot{\wEx}\big]\]
\[\dot{\mathbf J}=\wEth\begin{bmatrix}-s\theta&-c\theta&0\\c\theta&-s\theta&0\\0&0&0\end{bmatrix},\ \mathbf J^{-1}\dot{\mathbf J}=\wEth\begin{bmatrix}0&-1&0\\1&0&0\\0&0&0\end{bmatrix}\]
\end{fbx}
\begin{fbx}[Loads, power invariance \pg{32}]
$\wEQO^\top\dot{\wEx}=\wEQC^\top\wEz\Rightarrow\wEQO^\top\mathbf J\wEz=\wEQC^\top\wEz\ \forall\wEz\Rightarrow\boxed{\wEQC=\mathbf J^\top\wEQO}$, so $\wEQO=\mathbf J\wEQC$ (rotation: $\mathbf J^\top=\mathbf J^{-1}$).
\end{fbx}
\begin{fbx}[EoM in $\{O\}$ \pg{34}]
\[\boxed{{}^O\mathbf M\ddot{\wEx}+{}^Oc={}^O\mathbf Q},\ \ {}^O\mathbf M=(\mathbf J^\top)^{-1}\mathbf M\mathbf J^{-1},\ \ {}^Oc=(\mathbf J^\top)^{-1}\Big[c-\mathbf M\mathbf J^{-1}\dot{\mathbf J}\mathbf J^{-1}\dot{\wEx}\Big]\]
$C=G$: $c=\mathbf M[\boldsymbol\omega\times]\wEz$, so ${}^O\mathbf M=\operatorname{diag}(m,m,\wEI)$, ${}^Oc=\mathbf 0$: $m\ddot x=f_x$, $m\ddot y=f_y$, $\wEI\wEthh=\tau_z$. The $c$ in $\{C\}$ exists only because the axes turn. \chk
\end{fbx}
\begin{slideerr}[p.\,30 and p.\,34]
p.\,30 ``$\mathbf J(\eta)$'' should be $\mathbf J(\wEx)$. p.\,34 prints ${}^Oc$ without the $\mathbf J^{-1}$ in front of $\dot{\mathbf J}$ (test $C{=}G$, $m{=}10$, $\theta{=}90^\circ$, $u{=}1$, $v{=}0$, $\wEth{=}2$: correct $\mathbf 0$, printed $(-20,20,0)$). \chk
\end{slideerr}
\begin{worked}[loads between frames, ours]
$\theta{=}120^\circ$, $\wEQO{=}(4,-2,0.7)$: $\wEQC{=}\mathbf J^\top\wEQO{=}(\mathbf{-3.732},\mathbf{-2.464},0.7)$, $|\mathbf f|{=}4.472$ in both ($\mathbf J$ instead: wrong). \chk
\end{worked}

%==================================================================
\vbox{\hd{Euler integration step \pg{35}}
\begin{fbx}[Slide recipe \pg{35}]
\[\dot{\wEz}=\mathbf M^{-1}\big[\wEQC-c(\wEz)\big],\ \ \wEz\leftarrow\wEz+\Delta t\,\dot{\wEz},\ \ \dot{\wEx}=\mathbf J(\wEx)\wEz,\ \ \wEx\leftarrow\wEx+\Delta t\,\dot{\wEx}\]
\end{fbx}}
\begin{recipe}[One Euler step]
\begin{enumerate}
\item $c(\wEz_k)$ ($C=G$: $(-m\wEth v,\,m\wEth u,\,0)$); $\wEQC$ from wheels.
\item No slip: add the lateral reaction so $\dot v=0$.
\item $\dot{\wEz}=\mathbf M^{-1}(\wEQC-c)$; $\wEz_{k+1}=\wEz_k+\Delta t\dot{\wEz}$.
\item $\dot{\wEx}=\mathbf J(\theta_k)\wEz_{k+1}$; $\wEx_{k+1}=\wEx_k+\Delta t\dot{\wEx}$.
\item \textbf{Check:} radians; $v=0$ stays for diff drive.
\end{enumerate}
\end{recipe}
\begin{worked}[one step, ours]
$C{=}G$, $m{=}10$, $\wEI{=}0.25$, $r{=}0.1$, $a{=}0.2$; $\wEx{=}(1,2,\pi/6)$, $\wEz{=}(1,0,1.6)$; $\tau_L{=}0.4$, $\tau_R{=}0.6$; $\Delta t{=}0.1$.
$F{=}(4,6)$, $\wEQC{=}(10,0,0.4)$, $c{=}(0,16,0)$. Free body $\dot{\wEz}{=}(1,-1.6,1.6)$; no slip (${}^Cf_y{=}16$): $\dot{\wEz}{=}(\mathbf1,\mathbf0,\mathbf{1.6})$. $\wEz_1{=}(1.1,0,1.76)$; $\dot{\wEx}{=}(1.1c30^\circ,\,1.1s30^\circ,\,1.76){=}(0.9526,0.55,1.76)$; $\wEx_1{=}(\mathbf{1.0953},\mathbf{2.055},\mathbf{0.6996})$. (Old $\wEz_k$ in $\dot{\wEx}$: $(1.0866,2.05,0.6836)$; state your order.) \chk
\end{worked}

%==================================================================
\hd{Differential drive: wheel torques to chassis load \pg{36--37}}\nopagebreak[4]
\begin{fbx}[Differential drive \pg{37}]
\[\begin{bmatrix}{}^Cf_x\\{}^Cf_y\\{}^Cn_z\end{bmatrix}=\begin{bmatrix}1&1\\0&0\\-a&a\end{bmatrix}\begin{bmatrix}F_1\\F_2\end{bmatrix},\ \ F_i=\frac{\tau_i}{r}\]
\[\dot u=\frac{F_1+F_2}{m},\ \ \wEthh=\frac{a(F_2-F_1)}{\wEI}\ \ (C{=}G,\ \text{rest})\]
$F_1$ left wheel at $(0,+a)$, $F_2$ right at $(0,-a)$, both along $+\hat x$; $2a$ = track; column $=(d_x,d_y,x_id_y-y_id_x)$. Row 2 is zero: the sideways force is a no-slip reaction (not stated in the deck; say ``no slip $\Rightarrow v{=}\dot v{=}0$''). Power: $\boldsymbol\tau\cdot\dot{\boldsymbol\varphi}=\wEQC\cdot\wEz$.
\end{fbx}
\begin{recipe}[Torques $\to\dot u,\wEthh$]
\begin{enumerate}
\item $F_i=\tau_i/r$. \ 2.\ $f_x=F_1+F_2$, $n_z=a(F_2-F_1)$.
\item $\dot u=f_x/m$, $\wEthh=n_z/\wEI$, $v=0$; $R=\wEI f_x/(mn_z)$. \textbf{Check:} right wheel harder $\Rightarrow$ left turn.

\end{enumerate}
\end{recipe}
\begin{worked}[ours]
$m{=}10$, $\wEI{=}0.25$, $r{=}0.1$, $a{=}0.2$, $\tau_L{=}0.4$, $\tau_R{=}0.6$: $F{=}(4,6)$, $f_x{=}10$, $n_z{=}0.4$; $\dot u{=}\mathbf{1.0}$, $\wEthh{=}\mathbf{1.6}$, $R{=}0.625$. At $t{=}1$: $\dot\varphi_{L,R}{=}(1\mp0.32)/0.1{=}6.8,\,13.2$ rad/s; power $0.4(6.8){+}0.6(13.2){=}10.64\,\mathrm W{=}10(1){+}0.4(1.6)$. Holding $v{=}0$ needs $mu\wEth{=}16$ N sideways. \chk
\end{worked}
\begin{trap}
$2a$ as lever arm (doubles $\wEthh$); wrong sign of the $n_z$ row; forgetting $F=\tau/r$; $\mathbf J$ instead of $\mathbf J^\top$ for loads; calling a simulation from given forces ``inverse'' dynamics; $M_{33}=\wEI$ for an offset $G$.
\end{trap}

%==================================================================
\hd{Scenario: force to drive a circle \added}\nopagebreak[4]
\begin{recipe}[Circle of radius $R$ at speed $u$]
\begin{enumerate}
\item $\wEth=u/R$, $v=0$, $\dot u=\wEthh=0$.
\item $\wEQC=c=(0,\ mu\wEth,\ 0)=(0,\ mu^2/R,\ 0)$: force $mu^2/R$ sideways, towards the centre; supplied by the no-slip reaction, not the drive ($f_x=n_z=0$).
\item $\dot\varphi_{L,R}=(u\mp a\wEth)/r$.
\item \textbf{Check:} $R\to\infty$ gives $0$; double speed gives $4\times$ force.
\end{enumerate}
\end{recipe}
\begin{worked}[ours]
$m{=}12$, $u{=}0.6$, $R{=}1.5$: $\wEth{=}0.4$, $f_y{=}\mathbf{2.88}$ N. $a{=}0.2$, $r{=}0.1$: $\dot\varphi_L{=}\mathbf{5.2}$, $\dot\varphi_R{=}\mathbf{6.8}$ rad/s. \chk
\end{worked}

%==================================================================
\vbox{\hd{Exam-style questions}
\textbf{W13 sample Q4} (derive EoM), \textbf{Q10} (wheels): solved on p.\,\pageref{p:SAMPLE}. Marks not published; ~1 per step.}
\begin{examq}[1: offset COM \mk{~6 steps}]
$m{=}8$ kg, $\wEI{=}0.3$ kg\,m$^2$, $G$ at $({}^Cx_G,{}^Cy_G){=}(0.05,0.02)$ m. At $C$: ${}^C\mathbf f{=}(6,0)$ N, ${}^Cn_z{=}0.5$ N\,m, at rest. (a) Derive $\mathbf M$. (b) Find $\wEthh$, $\dot v$.
\solution
(a) At rest $\mathbf a_G=(\dot u-y_G\wEthh,\ \dot v+x_G\wEthh)$, $c=\mathbf 0$. Newton: $f_x=m(\dot u-y_G\wEthh)$, $f_y=m(\dot v+x_G\wEthh)$; Euler: $n_z=\wEI\wEthh-y_Gf_x+x_Gf_y$. Substitute:
$\mathbf M=\begin{bmatrix}m&0&-my_G\\0&m&mx_G\\-my_G&mx_G&\wEI+m(x_G^2+y_G^2)\end{bmatrix}=\begin{bmatrix}8&0&-0.16\\0&8&0.4\\-0.16&0.4&0.3232\end{bmatrix}$.
(b) $\dot v=-0.05\wEthh$, $\dot u=0.75+0.02\wEthh$; row 3: $-0.12+0.3\wEthh=0.5$.
\[\boxed{\wEthh=2.067\ \text{rad/s}^2,\ \ \dot v=-0.103,\ \ \dot u=0.791\ \text{m/s}^2}\]
Check: $\wEthh$ coefficient is exactly $\wEI=0.3$ after elimination. \chk
\end{examq}
\end{multicols}
